\section{Introduction}

To obtain physical observables in lattice field theory, it is essential to approximate the path integral as a sum over field configurations. Monte Carlo methods rely on random sampling from the probability distribution, determined by the action of the system~\cite{Knechtli:2017sna}. These methods can have high computational costs, particularly in the vicinity of a critical point in parameter space~\cite{Wolff:1989wq}. Generative models, as a class of machine learning (ML) algorithms, can be trained to generate new data that follow the same distribution as a given data set, representing the underlying target distribution~\cite{tomczak:2022deep}.
An important potential application of generative models in lattice QCD is therefore to improve the efficiency of Monte Carlo simulations~\cite{Boyda:2022nmh}.

Depending on the way the likelihood is estimated, roughly two main categories of generative models are being developed. In implicit maximum likelihood estimation (MLE), Generative Adversarial Networks (GANs) can learn to generate new configurations via a min-max game using training on a dataset prepared using, e.g., a standard lattice simulation.
A well-trained GAN can subsequently generate new samples, which should follow the same statistical distribution as the original dataset, but with reduced computational cost. Additionally, GANs can be used to improve the interpretability of lattice QCD simulations by generating visualizations of the system's behavior~\cite{Zhou:2018ill,Pawlowski:2018qxs}.
In explicit MLE, flow-based models have recently been proposed to improve the efficiency of lattice simulations~\cite{ Albergo:2019eim,  Kanwar:2020xzo, Nicoli:2020njz,Albergo:2021bna, Albergo:2021vyo, DelDebbio:2021qwf, Hackett:2021idh, Albergo:2022qfi, Bacchio:2022vje, Caselle:2022acb, Chen:2022ytr, Gerdes:2022eve, Albandea:2023wgd, Singha:2023cql}. Flow-based models are able to directly approach the underlying distribution of the quantum field without preparing training data. Although the method may encounter ``mode-collapse" and scalability problems~\cite{Abbott:2022zsh,Nicoli:2023qsl}, it allows for the generation of independent new configurations that can be used in MC simulations, again reducing the computational cost. The current progress in applying generative models also includes designing novel neural network architectures for specific field theories, e.g., autoregressive networks~\cite{Wang:2020hji,Luo:2023opo} and gauge-equivariant neural networks~\cite{Favoni:2020reg, Kanwar:2020xzo, Aronsson:2023rli, Abbott:2022zhs, Gerdes:2022eve,Luo:2023opo}. These approaches suggest that generative models have the potential to advance the capabilities of lattice QCD simulations. A comprehensive review can be found in Refs.~\cite{Zhou:2023pti, He:2023zin}.

Recently, the ML community has developed a new class of deep generative models known as Diffusion Models (DMs)~\cite{yang:2022diffusion}. Impressive success has been obtained in generating high-quality images via stochastic denoising processes, as showcased in prominent applications~\cite{Croitoru:2023dm} such as DALL-E 2~\cite{2022arXiv220406125R} and Stable Diffusion~\cite{Rombach_2022_CVPR}. As a prevalent class of generative models encoded in a stochastic process, DMs are grounded in Markov chains, and their implementation involves the utilization of variational inference methods. For an adoption in high-energy physics experiments, see, e.g., Refs.~\cite{Mikuni:2022xry,Mikuni:2023dvk}.


In numerical lattice field theory, there is an approach alternative to Monte Carlo methods to generate field configurations which relies on solving a stochastic differential equation, i.e., a Langevin equation. Its origin is stochastic quantization (SQ) as proposed by Parisi and Wu~\cite{Parisi:1980ys}. The basic idea of SQ is to view the quantum system as the thermal equilibrium limit of a hypothetical stochastic process with respect to a fictitious time. For gauge fields, it does not require gauge fixing~\cite{Parisi:1980ys}. Comprehensive reviews include Refs.~\cite{Damgaard:1987rr,Namiki:1993fd}.
For theories with a sign or complex weight problem, such as Quantum Chromodynamics at nonzero baryon density or theories in real (rather than Euclidean) time, the Langevin process can be extended to complex Langevin dynamics~\cite{Parisi:1983mgm}, which evades the sign problem in certain cases~\cite{Berges:2006xc,Aarts:2008rr,Aarts:2008wh,Seiler:2012wz,Sexty:2013ica}, see also the reviews~\cite{Aarts:2015tyj,Attanasio:2020spv,Berger:2019odf,Nagata:2021ugx}. Although the method has been put on firmer theoretical grounds~\cite{Aarts:2009uq,Aarts:2011ax,Nagata:2016vkn,Aarts:2017vrv,Scherzer:2018hid},
convergence and boundary condition issues continue to hinder obtaining a fully satisfying approach. Recent work to address this includes
Refs.~\cite{WesthHansen:2022iqd,Alvestad:2021hsi,Alvestad:2022abf,Lampl:2023xpb}.

In this work, we first point out the correspondence between generative DMs from the ML community and SQ in quantum field theory. Then, we explore using DMs to generate lattice field configurations. Specifically, we demonstrate for a two-dimensional $\phi^4$ lattice field theory that, combined with a hybrid Monte Carlo (HMC) algorithm, DMs can serve as an efficient global sampler along a Markov chain, leading to shorter autocorrelation times. We provide self-consistency checks and numerical evidence that this approach is %statistically exact and
able to capture the SQ perspective-induced stochastic dynamics of quantum field theories.

The paper is organized as follows. In the subsequent section, we provide a brief overview of the concept of SQ. In Sec.~\ref{sec:dm}, we present details of DMs. It is argued that the denoising process of DMs can explicitly serve as a SQ procedure. In Sec.~\ref{sec:phi4}, we explore the application of DMs to a two-dimensional scalar $\phi^4$ field theory.  Additionally, we demonstrate that DMs can accurately estimate the physical action in an unsupervised manner using the probabilistic flow approach. In Sec.~\ref{sec:sum}, we summarize the correspondence between DMs and SQ, and outline how this novel method can motivate further research, in particular for theories where it is expensive to sample from the prior distribution using standard approaches. App.~\ref{app:unet} contains information on the U-Net architecture employed, while Appendix \ref{app:she} contains a brief discussion on the Skilling-Hutchinson estimator.
A study of the dependence of observables and the acceptance rate on various system and DM parameters can be found in Appendix \ref{app:sta} and \ref{app:acc} respectively.
